\chapter*{Appendix 07: Coppice \& Carbon Commons}
\addcontentsline{toc}{chapter}{Appendix 07: Coppice \& Carbon Commons}

\section*{1. The Legacy Paradigm \& Its Failures}
The legacy timber industry uses linear extraction, clear-cutting forests and relying on diesel freight. Heating relies on fragile fossil-fuel grids. It extracts without accounting for the true Ecological Replacement Cost (ERC), leading to systemic biome collapse.

\section*{2. The Incremental Automation Pathway}
\textbf{Phase 1 (Manual Trust):} Communities manually track tree growth and agree verbally on how much to harvest, trusting each other not to over-cut.
\textbf{Phase 2 (AI-Assisted Policy):} Drone photogrammetry tracks canopy growth, and a human committee reviews the data to set seasonal harvest quotas.
\textbf{Phase 3 (Autonomous Cryptographic Enforcement):} The Digital Twin directly anchors Layer 5 policy. The smart contract mathematically rejects any \texttt{HarvestIntent} that exceeds the verified net biomass growth (ERC Hard Stop) without any human committee vote.

\section*{3. The Human Boundary (Limits of Automation)}
Sensors can track canopy mass and bio-acoustics, but they cannot fell a tree safely. Forestry labor requires human physical presence, intuition, and adherence to safety protocols. The system issues \texttt{HarvestIntent} bounties but relies on humans holding \texttt{Chainsaw\_Safety\_L1} credentials to execute the physical layer.

\section*{4. Orchestration Mechanics (The ``How'')}
\begin{itemize}
    \item \textbf{ERC Policy Enforcement:} The policy engine cross-references the requested biomass volume against the ledger's cryptographically minted growth tokens. This requires bridging L2 physical telemetry (drone scans) securely into L5 policy.
    \item \textbf{Credential Gating:} Execution requires a multi-signature transaction where one signature must come from a DID holding the specific \texttt{Forestry\_L2} Soulbound Token.
\end{itemize}
